\section*{Bab \ref{ch:b2:limb-organogenesis} --- Organogenesis: Membangun Tungkai Tetrapoda}

\begin{solution}{exo:b2:limb-organogenesis:1}
Proksimo-distal: \omterm{def:b2:limb-organogenesis:bud}{rabung ektoderm apikal} di sepanjang ujungnya, FGF.
Antero-posterior: \omterm{def:b2:limb-organogenesis:bud}{zona aktivitas pengutuban} pada tepi belakangnya,
\omterm{prop:b2:limb-organogenesis:zpa}{Sonic hedgehog}. Dorso-ventral: ektoderm punggungnya, Wnt7a.
\end{solution}

\begin{solution}{exo:b2:limb-organogenesis:2}
Stilopod: humerus, femur (Hoxa/d 9--10). Zeugopod: radius dan ulna,
tibia dan fibula (Hox 11). Autopod: pergelangan dan tangan, pergelangan
kaki dan kaki (Hox 13).
\end{solution}

\begin{solution}{exo:b2:limb-organogenesis:3}
Medan tungkainya ditetapkan (Tbx5); tunasnya tumbuh keluar di bawah FGF
rabungnya; ZPA dan ektoderm punggungnya memolakannya; sel yang
meninggalkan \omterm{prop:b2:limb-organogenesis:aer}{zona kemajuannya} memadat menjadi jari-jari; rawannya
berdiferensiasi; sendinya terbentuk di tempat rawannya ditekan;
pembuluh darahnya menyerbu dan tulang menggantikan rawannya; sel
antarjarinya mati; bakal otot dari somitnya bermigrasi masuk, melebur,
lalu melekat lewat tendon.
\end{solution}

\begin{solution}{exo:b2:limb-organogenesis:4}
(a) Sebuah tunggul dengan sedikit atau tanpa humerus; (b) humerus,
radius, dan ulna tetapi tanpa tangan; (c) sebuah tungkai yang lengkap
--- FGF menggantikan rabungnya.
\end{solution}

\begin{solution}{exo:b2:limb-organogenesis:5}
$x = 40\ln 2 = \qty{28}{\micro m}$, $40\ln 5 = \qty{64}{\micro m}$,
$40\ln 20 = \qty{120}{\micro m}$: wilayah selebar 28, 36, dan
\qty{56}{\micro m}.
\end{solution}

\begin{solution}{exo:b2:limb-organogenesis:6}
(a) 4-3-2-2-3-4, yaitu sebuah perangkat cermin yang penuh. (b) Sumber
yang lebih lemah memberi puncak yang lebih rendah, sehingga hanya jari
berambang rendah yang ditambahkan: 2-2-3-4. (c) Paparan yang pendek
juga hanya menetapkan jari tambahan yang paling depan: 2-2-3-4 ---
selnya mengintegralkan kepekatannya sepanjang waktu.
\end{solution}

\begin{solution}{exo:b2:limb-organogenesis:7}
$5000\times 2^{8} = 1.3\times 10^{6}$. Mencapai $4\times 10^{6}$
menuntut $\log_2 800 = 9.6$ kali pelipatduaan, sehingga daurnya
haruslah rata-rata kira-kira \qty{10}{h}, atau cacahnya haruslah
mencakup bakal otot yang bermigrasi masuk dari somitnya; kematian sel
akan memperburuk ketaksepadanannya, bukan memperbaikinya.
\end{solution}

\begin{solution}{exo:b2:limb-organogenesis:8}
Sel antarjari anak ayam menerima sebuah isyarat BMP lalu mati; sel
antarjari bebek mengungkapkan sebuah penghambat BMP lalu bertahan
hidup, sehingga selaputnya menetap. Penghambat pada sebuah manik
menyelamatkan selaput anak ayamnya: yaitu sebuah kaki berselaput.
Kematian sesuai jadwal pada sebuah cangkokan berarti selnya sudah
diperintah sebelum dibuang --- keputusannya dibuat dini lalu dijalankan
kemudian, secara otonom.
\end{solution}

\begin{solution}{exo:b2:limb-organogenesis:9}
Tanpa Hoxa13 dan Hoxd13: tanpa autopod --- tungkai depannya berakhir
pada pergelangannya. Tanpa Hoxa11 dan Hoxd11: radius dan ulna yang jauh
memendek dengan sebuah tangan yang melekat nyaris langsung pada
humerusnya. Tiap ruasnya membutuhkan pasangan Hox-nya sendiri;
daerahnya bukan sekadar penanda melainkan perintah.
\end{solution}

\begin{solution}{exo:b2:limb-organogenesis:10}
FGF dari rabungnya menjaga Shh tetap menyala di dalam ZPA-nya; Shh,
lewat sebuah penyambung mesoderm, menjaga FGF tetap menyala di dalam
rabungnya. Maka isyarat tiap pusatnya menyatakan bahwa yang lain sedang
bekerja, sehingga sebuah tungkai tidak pernah tumbuh tanpa dipolakan
maupun dipolakan tanpa tumbuh, dan kehilangan sementara salah satunya
diperbaiki yang lain. Kalau Shh-nya tidak dapat mempertahankan FGF-nya,
rabungnya akan menyusut dalam sehari dan pertumbuhan keluarnya akan
berhenti pada ruas mana pun yang telah tercapai: sebuah tungkai yang
terpenggal. Gelung persalinannya diakhiri peristiwa yang
digerakkannya --- kelahirannya membuang rangsangannya; gelung
tungkainya diakhiri diferensiasinya --- mesodermnya berhenti menyambung
ketika tungkainya lengkap.
\end{solution}

\begin{solution}{exo:b2:limb-organogenesis:11}
$e^{4r} = 8$: $r = \ln 8/4 = \qty{0.52}{d^{-1}}$. Pertumbuhannya:
pembiakan rawan jarinya yang lebih lama dan lebih cepat (sebuah penguat
tungkai yang menaikkan pengungkapan sebuah gen pertumbuhan di dalam
jarinya). Kematiannya: bertahan hidupnya jaringan antarjarinya lewat
pengungkapan sebuah penghambat BMP, yang mempertahankan selaputnya.
\end{solution}

\begin{solution}{exo:b2:limb-organogenesis:12}
Isyaratnya dipakai ulang karena sebuah modul reseptor--transduksi yang
bekerja murah untuk dipasang ulang: yang berubah antarjaringan bukanlah
isyaratnya melainkan perangkat gen yang selnya yang menerima berkompeten
menyalakannya, sehingga tampang kepekatan yang sama berarti neuron
motorik di satu tempat dan jari 4 di tempat lain. Maka sebuah \omterm{def:b2:genome-diversification:mutation}{mutasi}
pada gen semacam itu mengenai banyak organ sekaligus --- \omterm{def:b2:genome-diversification:mutation}{mutasi} Shh
memberi holoprosensefali dan cacat tungkai bersama-sama --- kecuali
kalau ia terletak di dalam sebuah penguat yang khusus jaringan, dan
dalam hal itu ia hanya mengubah satu organ, dan begitulah sebagian
besar perubahan evolusi pada gen ini terjadi.
\end{solution}

\begin{solution}{pb:b2:limb-organogenesis:1}
\textbf{1.} $96/12 = 8$ kali pelipatduaan: $2000\times 256 = 5.1\times 10^{5}$ sel.
\textbf{2.} $1000/256 \approx 4$: pembesaran sel dan matriks rawannya.
\textbf{3.} Panjangnya $\times 10$; sebuah silinder berjari-jari tetap
akan memberi $\times 10$ pada volumenya; kelebihan seratus kali
lipatnya berarti jari-jarinya juga tumbuh kira-kira sepuluh kali lipat
--- tunasnya melebar menjadi sebuah dayung sembari memanjang.
\textbf{4.} Pada 3 hari $300/400 = \qty{75}{\%}$; pada 7 hari
$300/4000 =
\qty{7.5}{\%}$.
\textbf{5.} Hanya ujungnya yang berada di bawah isyaratnya; sembari
ujungnya maju, sel yang tertinggal di belakang meninggalkan zonanya
dalam urutan pembuatannya, sehingga struktur proksimalnya ditetapkan
lebih dahulu dan struktur distalnya paling akhir.
\textbf{6.} Zeugopodnya (radius dan ulna), yang ditetapkan antara 3,5
dan 4 hari.
\textbf{7.} Stilopodnya pada 3,5 hari, zeugopodnya pada 4,0 hari,
autopodnya pada 4,5 hari (ujung jarinya lebih lambat).
\textbf{8.} Setengah hari: satu daur sel.
\textbf{9.} Tiap daur yang dihabiskan di dalam \omterm{prop:b2:limb-organogenesis:aer}{zona kemajuannya}
memajukan sandi Hox-nya (9--10, lalu 11, lalu 13); sel yang pergi
sesudah satu, dua, atau tiga daur mengusung sandi stilopod, zeugopod,
atau autopodnya.
\textbf{10.} Sebuah sayap yang lengkap: FGF-lah yang disediakan
rabungnya.
\textbf{11.} Sebuah pertumbuhan keluar kedua dari permukaan
punggungnya: struktur distal yang terganda, yaitu sebuah tungkai
tambahan.
\textbf{12.} Mesodermnya berhenti menyambungkan Shh ke rabungnya,
pengungkapan Shh-nya berakhir, rabungnya menyusut, dan selnya
berdiferensiasi alih-alih membelah: gelung yang menopang
pertumbuhannya telah putus.
\textbf{13.} $k = \ln 2/1740 = \qty{4.0e-4}{s^{-1}}$; $\lambda =
\sqrt{10^{-12}/4\times 10^{-4}} = \qty{50}{\micro m}$.
\textbf{14.} $x = 50\ln(1/0.6) = 26$, $50\ln 4 = 69$, $50\ln 12.5 =
\qty{126}{\micro m}$: wilayah selebar 26, 43, dan \qty{57}{\micro m}.
\textbf{15.} $600 - 126 = \qty{474}{\micro m}$.
\textbf{16.} Tiap batasnya bergerak ke luar sejauh $50\ln 2 = \qty{35}{\micro
m}$: menjadi 61, 104, 161. Ketiga lebarnya adalah 61, 43, dan \qty{57}{\micro m}:
hanya wilayah jari 4 yang tumbuh, yang lain sekadar tergeser.
Polanya bergeser ke depan alih-alih memperoleh sebuah jari, dan
daerah depan tanpa jarinya menyusut dari 474 menjadi \qty{439}{\micro m}.
\textbf{17.} Tiap perangkatnya membutuhkan \qty{126}{\micro m}:
sedikitnya \qty{252}{\micro m}; tunas \qty{600}{\micro m} memiliki
ruang, dengan \qty{348}{\micro m} jaringan tanpa jari di tengahnya.
\textbf{18.} Pada \qty{200}{\micro m} kedua landaiannya bertindih lalu
menjumlah: sel tengahnya melihat cukup banyak bagi jari 3 dan wilayah
jari 2-nya lenyap --- 4-3-4 atau 4-3-3-4, yaitu jari yang lebih sedikit
daripada sebuah perangkat cermin yang penuh.
\textbf{19.} Seperempat selnya memberi $c_0/4$: tidak ada sel yang
mencapai ambang jari 4 atau jari 3 ($0.25 c_0$ pada tepinya sendiri),
dan ambang jari 2-nya terlewati sampai $50\ln(0.25/0.08) =
\qty{57}{\micro m}$: yaitu satu jari 2 tambahan.
\textbf{20.} $3\times 300\times 300\times 200 = 5.4\times 10^{7}$
\unit{\micro m^3}: yaitu $5.4\times 10^{4}$ sel; sepanjang
\qty{86400}{s}, yaitu $0.6$ sel tiap sekon.
\textbf{21.} Manik BMP di dalam selaput bebeknya: selaputnya mati dan
jari kakinya terpisah; penghambat di dalam selaput anak ayamnya:
selaputnya bertahan hidup, yaitu sebuah kaki berselaput.
\textbf{22.} $e^{0.5\times 4} = e^{2} = 7.4$. Perubahan keduanya adalah
bertahan hidupnya selaput antarjarinya lewat sebuah penghambat BMP.
\textbf{23.} Fase akhir pengungkapan Hox 13 yang menetapkan sebuah
autopod, dan bertahannya rabungnya: lipatan ikannya membuat jari-jari
sirip alih-alih jari. Sirip \emph{Tiktaalik} memuat tulang mirip
pergelangan --- sebuah sirip yang mulai menjadi sebuah tungkai.
\textbf{24.} $5.4\times 10^{4}/1000 = 54$: $\log_2 54 = 5.8$ kali
pelipatduaan, yaitu kira-kira \qty{69}{h}, tiga hari untuk membuat apa
yang dibuang dalam sehari.
\textbf{25.} Delapan kali pelipatduaan; batas jarinya pada 26, 69, dan
\qty{126}{\micro m}; sebuah penggandaan cermin membutuhkan sedikitnya
\qty{252}{\micro m}.
\end{solution}
